\documentclass[10pt,a4paper]{amsart}
\usepackage[margin=19mm]{geometry}
\usepackage{amsmath,amssymb,mathtools}
\usepackage[scr=rsfs,cal=euler]{mathalfa}
\usepackage{xfrac}
\usepackage[unicode,hidelinks,bookmarks=false]{hyperref}
\newcommand{\ie}{i.e.\ }
\newcommand{\cf}{cf.\ }
\newcommand{\resp}{respectively\ }
\newcommand{\Subsection}{\subsection}
\providecommand{\nobreakdash}{\nobreak\hbox{-}}
\setlength{\emergencystretch}{2em}
\multlinegap0pt
\usepackage{amssymb}
\newtheorem{theorem}{Theorem}
\def\n{\mathbf n}
\def\r{\mathbb R}
\title{Classification of the ruled surfaces that are critical points of the Dirichlet energy}
\author{Rafael L\'opez}
\date{Source: arXiv:2605.14384v1 (CC BY 4.0)}
\begin{document}
\maketitle

\noindent\emph{Source and licence.} Classification of the ruled surfaces that are critical points of the Dirichlet energy. Version arXiv:2605.14384v1 (CC BY 4.0), distributed by arXiv under the Creative Commons Attribution 4.0 International licence, http://creativecommons.org/licenses/by/4.0/, as recorded in the arXiv licence metadata for this exact version and shown on the abstract page https://arxiv.org/abs/2605.14384v1. Full licence notice: \texttt{licenses/arxiv-2605.14384-CC-BY-4.0.txt}.\par\smallskip

\noindent\textit{Editorial scope.} Counted unit: the article's theorem t2 (source0.tex lines 225--227). The article's complete original proof of it is delivered with it (source0.tex lines 229--272, one \texttt{proof} environment of 44 lines). No mathematics is written, repaired or re-derived: the statement and the proof are the article's own lines, in the article's own notation, and no retained line is substituted or re-flowed.  Retained uncounted as the source-local context the proof needs: lines 87--89; lines 204--210; lines 220--222.  Nothing was omitted from any retained span, so the package carries no omission record.  No retained line contains a citation, so the wrapper carries no bibliography and no bibliography entry.  No retained line contains a figure environment, an image-inclusion command or drawing code, so no image is delivered and the package declares no asset dependency.  Editorial formatting only, in addition: an AMS \texttt{amsart} preamble that restates the article's own declarations and macros used by the retained lines, namely the environments \texttt{theorem} on separate counters, together with the standard packages those lines need.  The retained environments print in this document as Theorem 1 in order of appearance, whereas the article prints the same items with the numbering of its own full document.  The complete native source of the article as distributed in the arXiv e-print for this exact version is bundled unchanged as \texttt{sources/200580/source0.tex}.
\begin{equation}\label{fu}
\mathcal{F}(\Sigma)=\int_\Sigma \left(\dfrac{1}{\nu_3}-\nu_3\right)\, d\Sigma,
\end{equation}

\begin{equation}\label{h3}
\begin{split}
\Lambda\nu_3^3(1-\nu_3^2) \sqrt{\mbox{det}(g_{ij})}&=
2 (c_{11}^2h_{11}+2c_{11}\,c_{12}h_{12} +c_{12}^2h_{22})\\
&+2\nu_3^2 ( c_{21}^2h_{11}+2c_{21}\,c_{22}h_{12} +c_{22}^2h_{22}).
\end{split}
\end{equation}

\begin{equation}\label{para2}
X\colon I\times\r\to\r^3,\quad X(s,t)=\alpha(s)+t\vec{w},
\end{equation}

\begin{theorem} \label{t2}
Let $\Sigma$ be a cylindrical surface. If $\Sigma$ is a stationary surface of the Dirichlet energy \eqref{fu}, then $\alpha$ is a straight-line and $\Sigma$ is a plane ($\Lambda$) or $\alpha$ is a parabola and $\Sigma$ is a parabolic cylinder ($\Lambda\not=0$).
\end{theorem}

\begin{proof}
A first trivial case is when $\Sigma$ is a plane ($\Lambda=0$). From now, we will be discarded this case.   Consider the parametrization \eqref{para2}. Let $\n$ be the unit normal of $\alpha$ defined by $\n=\alpha''/|\alpha''|$, where $\kappa=|\alpha''|$ is the curvature of $\alpha$. Consider the orientation on  $\alpha$ such that $\alpha'\times \vec{w}=\n$. Since $X_s=\alpha'$ and $X_t=\vec{w}$, then $\nu=X_s\times X_t=\n$. Let write $e_3=(0,0,1)$ in coordinates with respect to $\{\alpha',\vec{w},\n\}$, 
$$e_3=e_{11}\alpha'+e_{22} \vec{w}+e_{33}\n.$$
Notice that $e_{33}=\langle e_3,\nu\rangle=\nu_3\not=0$.  
Recall that we discard the case $\nu_3^2= 1$ identically (horizontal plane). 

We compute all terms of \eqref{h3}. It is immediate  
 \begin{equation*}
\begin{split}
E_1&=e_{11}\alpha'+e_{22}\vec{w},\\
E_2&=-e_{22}\alpha'+e_{11}\vec{w}.
\end{split}
\end{equation*}
Since $h_{11}=\kappa$, $h_{12}=h_{22}=0$ and $\nu_3=e_{33}$, equation \eqref{h3} is 
\begin{equation}\label{ec}
\Lambda e_{33}^3(1-e_{33}^2)=2\kappa (e_{11}^2+e_{22}^2e_{33}^2).
\end{equation}
Consider a positively oriented (constant) basis $\mathcal{B}=\{v_1,v_2,v_3\}$ of $\r^3$ such that $v_3=\vec{w}$. Since the curve $\alpha$ is included in the $(v_1,v_2)$-plane and it is parametrized by arc-length, then there is a smooth function $\theta=\theta(s)$ such that 
$$\alpha'=\cos\theta v_1+\sin\theta v_2.$$
Thus 
$$\n=\alpha'\times \vec{w}=\sin\theta v_1-\cos\theta v_2.$$
Moreover, $\kappa=\langle\alpha'',\n\rangle=-\theta'$. Let write $e_3$ in coordinates with respect to $\mathcal{B}$, 
\begin{equation}\label{cij}
e_3=(a_1,a_2,a_3)=(\cos\varphi\cos\phi,\cos\varphi\sin\phi,\sin\varphi)
\end{equation}
for some $\varphi,\phi\in\r$.  Then 
\begin{equation*}
\begin{split}
e_{11}&=a_1\cos\theta+a_2\sin\theta,\\
 e_{22}&=a_3,\\
 e_{33}&=a_1\sin\theta-a_2\cos\theta.
 \end{split}
 \end{equation*}
Thus \eqref{ec} becomes 
\begin{equation}\label{ec2}
\begin{split}&\frac{\Lambda}{2}(a_1\sin\theta-a_2\cos\theta)^3(1-(a_1\sin\theta-a_2\cos\theta)^2)\\
&=-2\theta'\left((a_1\cos\theta+a_2\sin\theta)^2+a_3^2(a_1\sin\theta-a_2\cos\theta)^2\right).
 \end{split}
 \end{equation}
 Taking into account the value of $a_i$ in \eqref{cij}, equation \eqref{ec2} is  
\begin{equation}\label{ec3}
\theta'=-\frac{\Lambda}{2}(\cos\varphi)^3(\cos(\theta(s)+\phi))^3.\end{equation}
 If $\Lambda=0$, then $\kappa=-\theta'=0$ and $\alpha$ is a straight-line and $\Sigma$ is a plane. This case was already considered. If $\Lambda\not=0$,  equation \eqref{ec3} is the expression of  the curvature of a parabola when it is parametrized by arc-length. This proves the result.
 \end{proof}

\end{document}
